% GENERATED FILE. Edit canonical lessons, then run npm run generate:books.
% canonical-source-hash: fnv1a64:4570633a3deba077
% canonical-lessons: FR-C313-retraite, FR-C313-retraite-personne, FR-C313-vieillesse, FR-C313-hospitaliser, FR-C313-deces

\chapter{La retraite}
\label{ch:fr-la-retraite}
\hlchaptermodality{313}

\begin{chapteropening}
\textbf{By the end of this chapter:} \emph{I can talk about later life: retirement, retired people, old age, going into hospital and a death.}

Complete the last lesson of chapter 313: I can talk about later life: retirement, retired people, old age, going into hospital and a death.
\end{chapteropening}

\section[la retraite]{la retraite — retirement; the pension}

\label{lesson:FR-C313-retraite}



\begin{quote}
Before the new one: say the French for the tenant, then the French for the lease.
\end{quote}

\subsection*{You'll want to know: la retraite}
\textbf{la retraite} — \textquotedblleft{}retirement; the pension\textquotedblright{}.

\begin{quote}
Mon père a pris sa retraite l'année dernière.
\end{quote}

\begin{quote}
\emph{My father retired last year.}
\end{quote}

\emph{Prendre sa retraite} — to retire; \emph{être à la retraite} — to be retired.

\begin{cousinweb}
\emph{retraite} is the old past participle of Old French \emph{retraire} \textquotedblleft{}to withdraw\textquotedblright{}, which is inherited from Latin \emph{retrahere} \textquotedblleft{}to draw back\textquotedblright{} (Lewis \& Short).
\end{cousinweb}

\begin{grammarlens}[title={What you've built}]
One more word for telling, explaining and describing at B1.
\end{grammarlens}

\subsection*{Your turn}
\begin{itemize}
\raggedright
  \item \emph{Say it:} \emph{la retraite}
  \item \emph{Say it:} \emph{la retraite}, once more
  \item \emph{Say it:} a sentence of your own with \emph{la retraite}
  \item \emph{Recall:} say the French for the tenant, then the French for the lease, then say \emph{la retraite} again
\end{itemize}

\begin{tcolorbox}[breakable,colback=teal!4,colframe=teal!35!black,title={Before you move on}]
What does la retraite mean? (\textquotedblleft{}Retirement; the pension\textquotedblright{}.) Say it once more.
\end{tcolorbox}
\section[le retraité]{le retraité — the retired person, the pensioner}

\label{lesson:FR-C313-retraite-personne}



\begin{quote}
Before the new one: say the French for the lease, then the French for retirement.
\end{quote}

\subsection*{You'll want to know: le retraité}
\textbf{le retraité} — \textquotedblleft{}the retired person, the pensioner\textquotedblright{}.

\begin{quote}
Les retraités se retrouvent au café le matin.
\end{quote}

\begin{quote}
\emph{The pensioners meet at the café in the morning.}
\end{quote}

Feminine: \emph{la retraitée}.

\begin{cousinweb}
\emph{retraité} is built on \emph{retraite} \textquotedblleft{}retirement\textquotedblright{}, from Old French \emph{retraire}, inherited from Latin \emph{retrahere} \textquotedblleft{}to draw back\textquotedblright{} (Lewis \& Short).
\end{cousinweb}

\begin{grammarlens}[title={What you've built}]
One more word for telling, explaining and describing at B1.
\end{grammarlens}

\subsection*{Your turn}
\begin{itemize}
\raggedright
  \item \emph{Say it:} \emph{le retraité}
  \item \emph{Say it:} \emph{le retraité}, once more
  \item \emph{Say it:} a sentence of your own with \emph{le retraité}
  \item \emph{Recall:} say the French for the lease, then the French for retirement, then say \emph{le retraité} again
\end{itemize}

\begin{tcolorbox}[breakable,colback=teal!4,colframe=teal!35!black,title={Before you move on}]
What does le retraité mean? (\textquotedblleft{}The retired person, the pensioner\textquotedblright{}.) Say it once more.
\end{tcolorbox}
\section[la vieillesse]{la vieillesse — old age}

\label{lesson:FR-C313-vieillesse}



\begin{quote}
Before the new one: say the French for retirement, then the French for the retired person.
\end{quote}

\subsection*{You'll want to know: la vieillesse}
\textbf{la vieillesse} — \textquotedblleft{}old age\textquotedblright{}.

\begin{quote}
Il a eu une vieillesse heureuse.
\end{quote}

\begin{quote}
\emph{He had a happy old age.}
\end{quote}

The opposite of \emph{la jeunesse}.

\begin{cousinweb}
\emph{vieillesse} is built on \emph{vieil}, \emph{vieux} \textquotedblleft{}old\textquotedblright{}, inherited from Late Latin \emph{veclus}, from Latin \emph{vetulus} \textquotedblleft{}little old\textquotedblright{} (Lewis \& Short).
\end{cousinweb}

\begin{grammarlens}[title={What you've built}]
One more word for telling, explaining and describing at B1.
\end{grammarlens}

\subsection*{Your turn}
\begin{itemize}
\raggedright
  \item \emph{Say it:} \emph{la vieillesse}
  \item \emph{Say it:} \emph{la vieillesse}, once more
  \item \emph{Say it:} a sentence of your own with \emph{la vieillesse}
  \item \emph{Recall:} say the French for retirement, then the French for the retired person, then say \emph{la vieillesse} again
\end{itemize}

\begin{tcolorbox}[breakable,colback=teal!4,colframe=teal!35!black,title={Before you move on}]
What does la vieillesse mean? (\textquotedblleft{}Old age\textquotedblright{}.) Say it once more.
\end{tcolorbox}
\section[hospitaliser]{hospitaliser — to take into hospital, to hospitalise}

\label{lesson:FR-C313-hospitaliser}



\begin{quote}
Before the new one: say the French for the retired person, then the French for old age.
\end{quote}

\subsection*{You'll want to know: hospitaliser}
\textbf{hospitaliser} — \textquotedblleft{}to take into hospital, to hospitalise\textquotedblright{}.

\begin{quote}
Ma grand-mère a été hospitalisée pendant une semaine.
\end{quote}

\begin{quote}
\emph{My grandmother was in hospital for a week.}
\end{quote}

Mostly in the passive: \emph{il a été hospitalisé}.

\begin{cousinweb}
\emph{hospitaliser} is built on \emph{hospital}, the old form of \emph{hôpital}, from Latin \emph{hospitālis} \textquotedblleft{}of a guest\textquotedblright{}, built on \emph{hospes} \textquotedblleft{}guest, host\textquotedblright{} (Lewis \& Short).
\end{cousinweb}

\begin{grammarlens}[title={What you've built}]
One more word for telling, explaining and describing at B1.
\end{grammarlens}

\subsection*{Your turn}
\begin{itemize}
\raggedright
  \item \emph{Say it:} \emph{hospitaliser}
  \item \emph{Say it:} \emph{hospitaliser}, once more
  \item \emph{Say it:} a sentence of your own with \emph{hospitaliser}
  \item \emph{Recall:} say the French for the retired person, then the French for old age, then say \emph{hospitaliser} again
\end{itemize}

\begin{tcolorbox}[breakable,colback=teal!4,colframe=teal!35!black,title={Before you move on}]
What does hospitaliser mean? (\textquotedblleft{}To take into hospital, to hospitalise\textquotedblright{}.) Say it once more.
\end{tcolorbox}
\section[le décès]{le décès — the death (formal)}

\label{lesson:FR-C313-deces}



\begin{quote}
Before the new one: say the French for old age, then the French for to take into hospital.
\end{quote}

\subsection*{You'll want to know: le décès}
\textbf{le décès} — \textquotedblleft{}the death (formal)\textquotedblright{}.

\begin{quote}
Le décès a eu lieu à l'hôpital.
\end{quote}

\begin{quote}
\emph{The death took place at the hospital.}
\end{quote}

The formal and official word; in daily speech, \emph{la mort}.

\begin{cousinweb}
\emph{décès} is inherited from Latin \emph{dēcessus} \textquotedblleft{}a going away, a death\textquotedblright{} (EtymDB), from \emph{dēcēdere} \textquotedblleft{}to go away\textquotedblright{}.
\end{cousinweb}

\begin{grammarlens}[title={What you've built}]
One more word for telling, explaining and describing at B1.
\end{grammarlens}

\subsection*{Your turn}
\begin{itemize}
\raggedright
  \item \emph{Say it:} \emph{le décès}
  \item \emph{Say it:} \emph{le décès}, once more
  \item \emph{Say it:} a sentence of your own with \emph{le décès}
  \item \emph{Recall:} say the French for old age, then the French for to take into hospital, then say \emph{le décès} again
\end{itemize}

\begin{tcolorbox}[breakable,colback=teal!4,colframe=teal!35!black,title={Before you move on}]
What does le décès mean? (\textquotedblleft{}The death (formal)\textquotedblright{}.) Say it once more.
\end{tcolorbox}
